To evaluate our framework, we conducted four controlled experiments with a total of 110 Discord users across different community types (gaming, study groups, and professional communities). Each experiment tested specific hypotheses derived from our framework's claims.

Experiment 1 (User Engagement Study) examined whether proactive behavior improves user engagement. We recruited 50 users and randomly assigned them to either a reactive-only control group or a proactive treatment group for a two-week period. The proactive group received periodic suggestions based on system monitoring, while the control group only received responses to direct messages. We measured response latency, engagement rate, session length, and user satisfaction.

Experiment 2 (Safety and Trust Evaluation) assessed whether approval workflows increase perceived safety and trust. Using a within-subjects design with 30 users, each participant experienced both conditions: one where the agent could execute tools automatically upon suggestion, and one where approval was required before execution. We presented users with various scenarios ranging from beneficial actions to potentially harmful ones, and measured safety perceptions, trust scores, and actual prevented harmful actions.

Experiment 3 (External Knowledge Integration Study) evaluated whether integrating with external reasoning engines improves answer quality and reduces hallucinations. We created a benchmark of 100 factual questions across multiple domains and tested them with both an internal-knowledge-only version and a version augmented with NotebookLM access. Expert fact-checkers evaluated responses for factual accuracy, while we measured hallucination rates by comparing responses against known facts. We also assessed relevance and user satisfaction.

Experiment 4 (Component Ablation Study) determined the individual contribution of each framework component to overall performance. We tested 30 users each on five system configurations: the full system, the system without proactive behavior, the system without approval workflows, the system without external reasoning integration, and a reactive-only baseline. Users interacted with each configuration for equivalent periods, and we measured overall performance scores, user satisfaction, and task success rates.

All experiments received ethical review and informed consent from participants. Participants were compensated for their time at rates above local minimum wage. We ensured that no personal data was collected beyond what was necessary for the experiments, and all data was stored securely.